\documentclass[10pt]{article}
\usepackage[T1]{fontenc}
\usepackage[utf8]{inputenc}
\usepackage[french]{babel}
\usepackage{amsmath,amssymb,mathrsfs,mathpazo}
\usepackage{geometry}
\geometry{a4paper,margin=2cm}
\usepackage{graphicx,array,tabularx,tabularray,colortbl,multirow,multicol,booktabs,enumitem,pifont,ulem,fancybox,diagbox,cancel}
\usepackage{pstricks,pstricks-add,pst-plot,pst-tree,pst-node,pst-text,pst-func,pst-math,pst-bspline,pst-3dplot,pst-eucl,pst-solides3d}
\let\pscalculate\undefined
\usepackage{pst-bezier}
\usepackage{tikz,pgfplots,tkz-tab}
\usetikzlibrary{arrows,arrows.meta,patterns,calc,positioning,trees,decorations.markings,shapes.geometric,shapes.misc,shapes.arrows,intersections,decorations.pathmorphing,decorations.pathreplacing,backgrounds,shadows,babel}
\pgfplotsset{compat=1.18}
\usepgfplotslibrary{fillbetween}
\usepackage[np]{numprint}
\usepackage{listings,makecell,pas-tableur,twoopt}
\let\barre\undefined
\usepackage[pstricks]{tablvar}
\definecolor{darkgreen}{RGB}{0,100,0}
\providecommand{\red}{\color{red}}
\providecommand{\blue}{\color{blue}}
\providecommand{\green}{\color{green}}
\providecommand{\black}{\color{black}}
\providecommand{\white}{\color{white}}
\providecommand{\gray}{\color{gray}}
\providecommand{\cyan}{\color{cyan}}
\providecommand{\magenta}{\color{magenta}}
\providecommand{\yellow}{\color{yellow}}
\providecommand{\violet}{\color{violet}}
\providecommand{\orange}{\color{orange}}
\providecommand{\euro}{\texteuro}
\providecommand{\EUR}{\texteuro}
\providecommand{\ieme}{\textsuperscript{e}}
\providecommand{\degres}{\ensuremath{^\circ}}
\providecommand{\degree}{\ensuremath{^\circ}}
\providecommand{\vv}[1]{\overrightarrow{#1}}
\providecommand{\Coord}[1]{#1}
\providecommand{\decosix}{}
\providecommand{\starredbullet}{\ensuremath{\bullet}}
\pagestyle{empty}
\setlength{\parindent}{0pt}
\usepackage{adjustbox}\DeclareUnicodeCharacter{2236}{:}\DeclareUnicodeCharacter{2190}{\ensuremath{\leftarrow}}
\providecommand{\R}{\mathbb{R}}
\providecommand{\N}{\mathbb{N}}
\providecommand{\Z}{\mathbb{Z}}
\providecommand{\Q}{\mathbb{Q}}
\providecommand{\D}{\mathbb{D}}
\providecommand{\C}{\mathbb{C}}
\providecommand{\e}{{\mathrm{e}}}
\providecommand{\vect}[1]{\overrightarrow{#1}}
\providecommand{\vectt}[1]{\overrightarrow{\mathrm{#1}}}
\providecommand{\vvt}[1]{\overrightarrow{\mathrm{#1}}}
\providecommand{\np}[2][]{{#2~#1}}
\providecommand{\vv}[1]{{\overrightarrow{#1}}}
\renewcommand{\d}{{\mathrm{d}}}
\providecommand{\degres}{{^\circ}}
\providecommand{\degre}{{^\circ}}
\providecommand{\blue}{\color{blue}}
\providecommand{\red}{\color{red}}
\providecommand{\enskip}{\quad}
\providecommand{\llbracket}{\lbrack\!\lbrack}
\providecommand{\rrbracket}{\rbrack\!\rbrack}
\providecommand{\euro}{{\text{€}}}
\providecommand{\no}{{\text{n°}}}
\providecommand{\interval}[2]{{[#1;#2]}}
\providecommand{\intervFO}[2]{{[#1;#2[}}
\providecommand{\dotfill}{\ldots}

\providecommand{\ds}{\displaystyle}
\providecommand{\textding}[1]{#1}
\providecommand{\Oij}{\(\mathrm{O};\vec{\imath},\vec{\jmath}\)}
\providecommand{\Oijk}{\(\mathrm{O};\vec{\imath},\vec{\jmath},\vec{k}\)}
\providecommand{\Ouv}{\(\mathrm{O};\vec{u},\vec{v}\)}
\providecommand{\intervFF}[2]{\ensuremath{[#1;#2]}}
\renewcommand{\euro}{\ifmmode\text{\texteuro}\else\texteuro\fi}
\DeclareUnicodeCharacter{00B0}{\ifmmode^{\circ}\else\textdegree\fi}
\DeclareUnicodeCharacter{20AC}{\ifmmode\text{\texteuro}\else\texteuro\fi}
% Énoncés : Éducation nationale. Transcription et corrigés : APMEP, Denis Vergès et contributeurs. Découpage et classement : CoopMaths, Sébastien Lozano, Éric Elter et contributeurs. Adaptation de la présentation : Maths Explorer.
% CC BY-SA : https://coopmaths.fr/www/about/
% https://coopmaths.fr/alea/static/bac/2022/tex/bac_2022_05_sujet1_polynesie_1.tex
\begin{document}
\noindent\textit{APMEP / CoopMaths. CC BY-SA. Présentation adaptée par Maths Explorer.}\par\medskip
\medskip

\emph{Cet exercice est un questionnaire à choix multiples. Pour chacune des six questions suivantes, une seule des quatre réponses proposées est exacte.\\
Une réponse fausse, une réponse multiple ou l'absence de réponse à une question ne rapporte ni n'enlève de point.\\
Pour répondre, indiquer sur la copie le numéro de la question et la lettre de la réponse choisie. Aucune justification n'est demandée.}

\medskip

\begin{enumerate}
\item On considère  la fonction $g$ définie et dérivable sur $]0~;~+ \infty[$ par :

\[g(x) = \ln \left(x^2 + x + 1\right).\]

Pour tout nombre réel $x$ strictement positif :

\begin{center}
\begin{tabularx}{\linewidth}{X X}
\textbf{a.~~}$g'(x) = \dfrac{1}{2x + 1}$&\textbf{b.~~}$g'(x) = \dfrac{1}{x^2 + x + 1}$\\[8pt]
\textbf{c.~~}$g'(x) = \ln (2x + 1)$&\textbf{d.~~}$g'(x) = \dfrac{2x + 1}{x^2 + x + 1}$
\end{tabularx}
\end{center}
\item La fonction $x \longmapsto \ln (x)$ admet pour primitive sur $]0~;~+ \infty[$ la fonction :

\begin{center}
\begin{tabularx}{\linewidth}{*{4}{X} }
\textbf{a.~~}$x \longmapsto \ln (x)$&\textbf{b.~~}$x \longmapsto \dfrac{1}{x}$&\textbf{c.~~}$x \longmapsto x \ln (x) - x$&\textbf{d.~~}$x \longmapsto \dfrac{\ln (x)}{x}$
\end{tabularx}
\end{center}
\item On considère la suite $\left(a_n\right)$ définie  pour tout $n$ dans $\N$ par :

\[a_n = \dfrac{1 - 3^n}{1 + 2^n}.\]

La limite de la suite $\left(a_n\right)$ est égale à :

\begin{center}
\begin{tabularx}{\linewidth}{*{4}{X} }
\textbf{a.~~}$- \infty$&\textbf{b.~~}$- 1$&\textbf{c.~~}$1$&\textbf{d.~~}$+ \infty$
\end{tabularx}
\end{center}
\item 
On considère une fonction $f$ définie et dérivable sur $[-2~;~2]$. Le tableau de variations de la fonction $f'$ dérivée de la fonction $f$ sur l'intervalle $[2~;~2]$ est donné par :

\begin{center}
{\renewcommand{\arraystretch}{1.2}
\psset{nodesep=3pt,arrowsize=2pt 3}  
\def\esp{\hspace*{1.5cm}}
\def\hauteur{0pt}
$\begin{array}{|c| *4{c} c|}
\hline
 x & -2 & \esp & 0 & \esp & 2 \\
\hline
  & \Rnode{max1}{1}  &  &  &  & \Rnode{max2}{-1}   \\
\text{variations de }f' & &  & & &  \rule{0pt}{\hauteur} \\
 &  & &   \Rnode{min}{-2} & & \rule{0pt}{\hauteur}
\ncline{->}{max1}{min} \ncline{->}{min}{max2}
\rput*(-3.7,0.6){\Rnode{zero}{0}}
\rput(-3.7,1.85){\Rnode{alpha}{-1}}
\\
\hline
\end{array}$
}
\end{center}


La fonction $f$ est :

\begin{center}
\begin{tabularx}{\linewidth}{X X}
\textbf{a.~~} convexe sur $[- 2~;~- 1]$&\textbf{b.~~} concave sur [0~;~1]\\
\textbf{c.~~} convexe sur $[- 1~;~2]$&\textbf{d.~~}concave sur $[-2~;~0]$
\end{tabularx}
\end{center}

\item On donne ci-dessus la courbe représentative de la dérivée $f'$ d'une fonction $f$ définie sur l'intervalle $[-2~;~4]$.

\begin{center}
\psset{unit=1cm,arrowsize=2pt 3}
\begin{pspicture*}(-2.2,-3.2)(4.2,3.2)
\psgrid[gridlabels=0pt,subgriddiv=1,gridcolor=lightgray]
\psaxes[linewidth=1.25pt,labelFontSize=\scriptstyle]{->}(0,0)(-2.2,-3.2)(4.2,3.2)
\psbcurve[plotpoints=5000,linewidth=1.25pt,linecolor=blue](-2,-3)(-1,0)(0,1)L(0.1,1)(1,0)L(1.5,-1)(2,-1)L(2.1,-1)(3,0)(4,3)

\end{pspicture*}
\end{center}

Par lecture graphique de la courbe de $f'$, déterminer l'affirmation correcte pour $f$ :

\begin{center}
\begin{tabularx}{\linewidth}{X X}
\textbf{a.~~} $f$ est décroissante sur [0~;~2]&\textbf{b.~~}$f$ est décroissante sur $[-1~;~0]$\\
\textbf{c.~~}$f$ admet un maximum en 1 sur [0~;2]&\textbf{d.~~}$f$ admet un maximum en 3 sur [2~;~4]
\end{tabularx}
\end{center}

\item Une action est cotée à 57 \euro. Sa valeur augmente de 3\,\% tous les mois.

La fonction python seuil() qui renvoie le nombre de mois à attendre pour que sa valeur dépasse 200 \euro{} est :

\begin{center}
\begin{tabularx}{\linewidth}{X X}
\textbf{a.~~}&\textbf{b.~~}\\
\fbox{\begin{tabular}{l}
def seuil() :\\
\quad m=0\\
\quad v=57\\
\quad while v < 200 :\\
\qquad m=m+1\\
\qquad v = v*1.03\\
\quad return m
\end{tabular}
} & \fbox{\begin{tabular}{l}
def seuil() :\\
\quad m=0\\
\quad v=57\\
\quad while v > 200 :\\
\qquad m=m+1\\
\qquad v = v*1.03\\
\quad return m
\end{tabular}
}  \\
\textbf{c.~~} &\textbf{d.~~} \\
\fbox{\begin{tabular}{l}
def seuil() :\\
\quad v=57\\
\quad for i in range (200) :\\
\qquad v = v*1.03\\
\quad return v
\end{tabular}
} &\fbox{\begin{tabular}{l}
def seuil() :\\
\quad m=0\\
\quad v=57\\
\quad if v < 200 :\\
\qquad m=m+1\\
\quad else :\\
\qquad v = v*1.03\\
\quad return m
\end{tabular}
}
\end{tabularx}
\end{center}
\end{enumerate}

\bigskip

\end{document}
